\documentclass[11pt]{article}

\usepackage[margin=1.15in]{geometry}
\usepackage{amsmath,amssymb,amsthm}
\usepackage{booktabs}
\usepackage{array}
\usepackage{enumitem}
\usepackage{microtype}
\usepackage{setspace}
\usepackage[round]{natbib}
\usepackage{xcolor}
\usepackage[colorlinks=true,linkcolor=blue!45!black,citecolor=blue!45!black,urlcolor=blue!45!black]{hyperref}

\onehalfspacing

\newtheorem{theorem}{Theorem}[section]
\newtheorem{proposition}[theorem]{Proposition}
\newtheorem{lemma}[theorem]{Lemma}
\newtheorem{corollary}[theorem]{Corollary}
\theoremstyle{definition}
\newtheorem{assumption}{Assumption}
\newtheorem{definition}[theorem]{Definition}
\newtheorem{example}[theorem]{Example}
\theoremstyle{remark}
\newtheorem{remark}[theorem]{Remark}
\newtheorem{observation}[theorem]{Observation}

\newcommand{\ubar}[1]{\underline{#1}}
\newcommand{\R}{\mathbb{R}}
\newcommand{\E}{\mathbb{E}}
\newcommand{\Kcal}{\mathcal{K}}
\newcommand{\dd}{\,d}
\newcommand{\ubc}{\ubar{c}}
\newcommand{\uba}{\ubar{a}}

\title{Equilibrium Transition Dynamics in Heterogeneous-Agent
Economies}
\author{Abdoulaye Ndiaye using Claude Fable 5}
\date{July 20, 2026}

\begin{document}
\maketitle

\begin{abstract}
\noindent
Quantitative work with heterogeneous-agent models routinely computes
perfect-foresight transition paths: an interest-rate path clears the
asset market at every date while the wealth distribution evolves from
given initial conditions. We ask when such equilibrium paths exist. In
the continuous-time Huggett economy with CRRA utility we show that
market clearing stated in \emph{stocks} is a degenerate restriction on
the rate path: today's interest rate has no contemporaneous effect on
today's aggregate wealth. The operative form of clearing is a
\emph{flow} condition---aggregate consumption equals aggregate income
plus interest on the bond supply---in which the rate's only
instantaneous leverage comes from the bond stock and from
hand-to-mouth households, whose interest burden moves one-for-one with
the rate. We prove that equilibrium transitions exist and are unique
over short horizons whenever the debt of hand-to-mouth households is
smaller than the bond supply, and that whether households remain at
the borrowing constraint near the end of a transition is governed by
terminal conditions rather than by discounting. With risk aversion
above one, monotonicity of market clearing in the rate path fails on
known economies, and we argue that multiple self-fulfilling transition
paths connecting the same endpoints become a live possibility.\\[0.5em]
\emph{JEL codes:} C62, D52, E21, D31.\quad
\emph{Keywords:} incomplete markets, borrowing constraints, transition
dynamics, existence of equilibrium, mean field games.
\end{abstract}

\newpage

%======================================================================
\section{Introduction}
\label{sec:intro}

Modern quantitative macroeconomics runs on transition paths. To
evaluate a credit crunch, a fiscal reform, or a monetary shock in a
Bewley--Huggett--Aiyagari economy, one computes the perfect-foresight
response: starting from an initial distribution of wealth, households
optimize against a conjectured path of prices, the wealth distribution
evolves under their decisions, and the price path is adjusted until
markets clear at every date \citep{KMV2018,BoppartKrusellMitman2018,
AuclertEtAl2021,AHLLM2022}. The algorithms that do this are by now
standard. The theory behind them is not: while existence and (under
conditions) uniqueness of \emph{stationary} equilibria are well
understood \citep{Aiyagari1994,Acikgoz2018,AHLLM2022,Hofer2025}, there
is, to our knowledge, no theorem guaranteeing that the equilibrium
transition path being computed exists, for any horizon, any preference
parameters, or any level of asset supply. This paper studies that
question in the canonical laboratory: the continuous-time Huggett
economy with CRRA utility, two income states, a borrowing constraint,
and a bond in fixed supply $B\ge0$, over a horizon $[0,T]$.

Why is this not a routine extension of the steady-state theory? In a
stationary equilibrium the entire influence of the wealth
distribution on any single household is summarized by one number, the
interest rate $r$. Fixing $r$ decouples the household problem from
the distribution, and existence reduces to finding a root of a scalar
excess-supply function---Aiyagari's celebrated diagram, formalized by
an intermediate-value argument. Along a transition, the unknown is a
\emph{path} $r(\cdot)$, pinned down by a clearing condition at a
continuum of dates. Three economic facts, which we formalize, make
this fixed-point problem behave very differently from its scalar
ancestor.

\emph{First, almost nobody responds to today's interest rate today.}
A household's consumption at date $t$ is determined by its wealth and
by the price path it expects \emph{from $t$ onward}; changing $r$ at
the single instant $t$ leaves its date-$t$ behavior unchanged. The
wealth distribution at $t$, in turn, depends on prices \emph{before}
$t$. Consequently today's rate has no contemporaneous effect on
today's aggregate wealth: market clearing stated in stocks
(``aggregate wealth equals bond supply at every date'') is a
degenerate restriction on the rate path---in technical language, a
first-kind integral equation, which admits no stable inversion
(Theorem~\ref{thm:firstkind}). The exception, and it is the crux of
the paper, is a group whose behavior \emph{does} depend on today's
rate: hand-to-mouth households at the borrowing constraint, who
consume their income flow $y_1+r(t)\uba$ and whose interest burden
therefore moves one-for-one with $r(t)$; and bondholders in the
aggregate, whose interest income moves with $r(t)B$.

\emph{Second, the operative form of market clearing is a flow
condition.} Differentiating aggregate wealth along the distributional
dynamics shows that clearing at every date is equivalent to clearing
at the initial date---a restriction on the \emph{data}, not on
prices---plus the aggregate budget identity
\[
C(t)\;=\;Y(t)+r(t)B \qquad\text{at every }t,
\]
aggregate consumption equals aggregate income plus interest on the
bond supply (Lemma~\ref{lem:flow}). In this form the rate has exactly
the two contemporaneous levers just described: the sensitivity of
excess demand to $r(t)$ is $B+|\uba|\,m_1(t)$, where $m_1(t)$ is the
mass of constrained households. All of our existence results are built
on this diagonal. It is telling that numerical practice discovered the
same structure by trial: the transition algorithm of
\citet[\S5.5]{AHLLM2022} updates the rate path using the \emph{time
derivative} of excess supply---the flow---rather than its level.

\emph{Third, the borrowing constraint itself becomes a moving part.}
In the stationary economy, impatience ($r<\rho$) guarantees that the
constraint binds and that a point mass of hand-to-mouth households
forms at the borrowing limit. Along a transition this is no longer
governed by discounting. We show that whether households stay at the
constraint near the end of the horizon is controlled by the
\emph{terminal condition}: with the standard choice---continuation
values from a stationary equilibrium at rate $r^*$---the constrained
group stays put at the end of the horizon if and only if
$r(T)\ge r^*$, and disperses whenever the path ends below $r^*$, even
if $r(t)<\rho$ throughout (Corollary~\ref{cor:knife}). The economics
is elementary once seen: constrained households are debtors, so a
rate \emph{below} the one embedded in their continuation values
lowers their debt service, leaves them income to spare, and they save
their way off the constraint. Since the constrained mass $m_1(t)$ is
also the rate's main lever on market clearing, the model's two
degeneracies---at the borrowing limit and in the price---are coupled:
episodes in which the constrained group disperses are episodes in
which the market loses its contemporaneous handle on the interest
rate.

\paragraph{Results.} With this structure in place, our main results
are as follows.

(1) \emph{Existence and uniqueness for short transitions with
positive bond supply} (Theorem~\ref{thm:smallT}). If the initial
distribution clears the market, if the data imply an interior
initial rate (an ``anchoring'' condition), and if
\[
\theta\;:=\;\frac{|\uba|\cdot\sup_t m_1(t)}{B}\;<\;1,
\]
then an equilibrium transition exists and is unique on $[0,T]$ for
$T$ below an explicit threshold. The condition $\theta<1$ has a
transparent reading: a rise $dr$ in the rate mechanically raises the
consumption the economy must absorb by $B\,dr$ (interest paid on the
bond supply) and lowers hand-to-mouth consumption by
$|\uba|m_1\,dr$ (interest charged on constrained debt); the
equilibrium feedback loop is stable when the second effect is smaller
than the first---equivalently, when the debt of hand-to-mouth
households is smaller than the bond supply. That is the stability
condition of the transition.

(2) \emph{Zero net supply lives or dies by the constrained mass}
(Theorem~\ref{thm:B0atom}, Proposition~\ref{prop:notheorem}). With
$B=0$, the flow condition pins aggregate consumption to aggregate
income, and the \emph{only} contemporaneous lever of the rate is the
constrained group. If the initial distribution carries a point mass
at the constraint and the constraint remains binding, that mass decays
no faster than the rate at which its members draw the high income
state, and a short-horizon equilibrium exists. If instead $B=0$ and no
one is constrained on some time interval, market clearing there
contains today's rate \emph{nowhere}: the constraint becomes a
restriction linking only past and future, no local solution theory
applies, and---consistent with the nonexistence results of
\citet{Ambrose2021} for a frictionless relative of this model---exact
clearing is generically the wrong demand. Relaxations (a slightly
elastic bond supply, or a price rule responding to excess demand)
restore a well-posed problem; we make precise the sense in which they
are a necessity rather than a computational convenience.

(3) \emph{The household problem along a price path is harder than its
stationary counterpart, in identifiable places}
(Section~\ref{sec:household}, Appendices~\ref{app:household} and
\ref{app:distribution}). We provide a complete well-posedness package:
the household's value function is the unique constrained viscosity
solution of its Bellman equation along any admissible rate path;
marginal values admit two-sided bounds, constructed by explicit
barrier arguments because the stationary shortcut---the binding
constraint pinning the boundary gradient---is unavailable; consumption
near the constraint follows a square-root law whose speed coefficient
$\nu_1(t)$ acquires a new term $\dot r(t)\uba$ (a \emph{rising} rate
slows the approach to the constraint through the debtor's income
channel); and the wealth distribution evolves as a measure with a
point mass at the constraint that grows while the constraint binds and
detaches \emph{discontinuously} when it stops binding. Two features of
the finite horizon are genuinely simpler and worth recording: neither
$\rho>0$ nor $r<\rho$ plays any role in household well-posedness.

(4) \emph{Long transitions bifurcate on risk aversion}
(Section~\ref{sec:long}). For $\gamma\le1$ (intertemporal elasticity
of substitution at least one---the range in which the stationary
equilibrium is known to be unique), a natural monotonicity property of
aggregate consumption in the rate path, if it holds, localizes all
equilibria between two computable paths and yields uniqueness under a
verifiable collapse condition (Theorem~\ref{thm:mono}); we explain
exactly why this monotonicity is a hypothesis rather than a theorem in
finite horizon. For $\gamma>1$ it is not available even as a
hypothesis: \citet{Hofer2025} constructs economies whose stationary
aggregate saving \emph{decreases} in the interest rate over a range,
and on long horizons such economies produce singular linearizations of
the market-clearing map and, plausibly, multiple self-fulfilling
transition paths connecting the same endpoints---anticipation-driven
loops in which a believed future path of rates alters consumption
today, reshapes the wealth distribution, and validates the belief
(Proposition~\ref{prop:multiplicity}). For risk aversion above
one---the empirically common calibration---uniqueness of transition
dynamics should be regarded as an assumption of computational
practice, not a theorem.

\subsection{Related literature}
\label{sec:literature}

The stationary theory is by now well developed.
\citet{AHLLM2022} formulate the continuous-time system, prove
existence of stationary equilibria by the Aiyagari intermediate-value
argument, and prove uniqueness when the intertemporal elasticity of
substitution is at least one and borrowing is prohibited;
\citet{Acikgoz2018} rigorizes the discrete-time counterpart and gives
a calibrated multiplicity example. Mathematically complete stationary
treatments are given by \citet{Shigeta2023} (bounded utility) and
\citet{Hofer2025} (CRRA utility under a no-borrowing limit; ergodicity
of the wealth process, continuity of aggregate assets in $r$, and the
non-monotonicity construction for $\gamma>1$ used in
Section~\ref{sec:long}); \citet{BayerRendallWaelde2019} analyze the
invariant wealth distribution at fixed prices. All of these concern
steady states.

For transitions, the closest prior work is \citet{Ambrose2021}, who
studies the household-wealth model of \citet{ABLLM2014}---a variant
with diffusive income, \emph{no} borrowing constraint (compactly
supported wealth densities), and zero net supply. He proves small-time
existence and uniqueness for a \emph{relaxed} problem in which exact
clearing is replaced by constancy of the flow imbalance, and gives a
nonexistence result (his \S8) for the exact problem: with generic
data, no exactly-clearing solution exists over short horizons. In our
language, his model has $B=0$ and, by construction, no mass at a
borrowing constraint, so the rate's contemporaneous lever on clearing
is identically zero; his Remark~1 notes explicitly that mass at a
wealth boundary would contribute ``another term proportional to
$r$'' through which clearing could be enforced. The present paper can
be read as the constrained-economy counterpart: the point mass of
hand-to-mouth households (and positive bond supply) restores exactly
that missing term, and with it an existence theory for exact
clearing. \citet{ChenEtAl2025} obtain a time-dependent
market-clearing rate path in an overlapping-generations economy, under
a never-binding natural borrowing limit and a short-lifespan
restriction that plays the role of our short horizon.

Methodologically the paper sits in the mean field games (MFG)
literature. Finite-horizon well-posedness results exist for
price-mediated MFG in which the price is a multiplier on a constraint
involving agents' \emph{controls}
\citep{GomesSaude2021,FujiiTakahashi2022}, for oligopoly models with
prices that are explicit functionals of the distribution
\citep{GraberBensoussan2018}, for MFG of controls
\citep{CardaliaguetLehalle2018,Kobeissi2022}, including a
state-constrained case under a smallness condition on the dependence
of preferences on the joint law of states and controls
\citep{GraberMayorga2021}---a condition that, in our multiplier
formulation, scales like $1/B$ and fails precisely in the
economically central small-$B$ regime--- and for general non-separable
forward--backward parabolic systems over short horizons
\citep{CirantGianniMannucci2020}. In all of these, either the
clearing constraint responds nondegenerately to the contemporaneous
price through the controls, or there is no exact clearing constraint
at all; none features the combination present here---a hard borrowing
constraint with a boundary point mass, switching income risk, a
Hamiltonian singular at zero marginal value, and a price path pinned
by instant-by-instant clearing in stocks. The deterministic theory of
MFG with state constraints
\citep{CannarsaCapuani2018,CCC2021} supplies the right language for
constrained trajectories; the planning problem
\citep{Porretta2014} treats a single terminal constraint enforced by a
multiplier, where our clearing imposes one scalar constraint per
instant.

Finally, our degeneracy theorem speaks to computational practice. In
discrete time, the sequence-space Jacobian of
\citet{AuclertEtAl2021} is precisely a discretization of the
linearized market-clearing map that our Theorem~\ref{thm:firstkind}
analyzes; the theorem implies that in the continuous-time limit the
Jacobian of \emph{stock}-level clearing has vanishing diagonal and is
compact---ill-conditioned at fine time discretizations---while the
flow form retains an invertible diagonal $B+|\uba|m_1(t)$. This is
consistent with, and explains, the practice of updating rate paths
via flow imbalances \citep[\S5.5]{AHLLM2022}.

The paper proceeds as follows. Section~\ref{sec:model} lays out the
economy. Section~\ref{sec:household} analyzes the household problem
along a rate path, with the binding/dispersal results.
Section~\ref{sec:clearing} develops the structure of market clearing.
Section~\ref{sec:smallT} contains the short-horizon existence and
uniqueness theory, Section~\ref{sec:long} the long-horizon analysis.
Section~\ref{sec:conclusion} concludes. Appendices~\ref{app:household}
--\ref{app:proofs} contain the technical package and proof
architectures.

%======================================================================
\section{The economy}
\label{sec:model}

\subsection{Households}

Time is continuous over a finite horizon $[0,T]$. There is a continuum
of households with CRRA preferences
\[
\E_0\int_0^\infty e^{-\rho t}\,u(c_t)\dd t,\qquad
u(c)=\frac{c^{1-\gamma}}{1-\gamma},\ \gamma>0
\quad(u=\log\text{ for }\gamma=1),\qquad \rho\ge0 .
\]
A household's income $y_t$ takes two values $0<y_1<y_2$
(``unemployed''/``employed''), switching with Poisson intensities
$\lambda_1,\lambda_2>0$. Wealth is held in a riskless bond and evolves
as
\[
\dot a_t=y_t+r(t)\,a_t-c_t,\qquad a_t\ \ge\ \uba,
\qquad -\infty<\uba<0,
\]
where $r(t)$ is the market interest rate, taken as given by the
household, and $\uba$ is an exogenous borrowing limit. The horizon is
closed by a terminal valuation $v_{T,j}(a)$ of end-of-horizon wealth;
the canonical choice, following \citet{AHLLM2022}, is the value
function $v_{\infty}(\cdot\,;r^*)$ of a stationary equilibrium with
rate $r^*$, so that the finite-horizon problem is a truncation of a
transition converging to a steady state.

The household's value functions $v_j(a,t)$, $j=1,2$, solve the
backward Bellman (HJB) system
\begin{equation}\label{eq:HJB}
\rho v_j=\partial_t v_j+\max_{c\ge0}\Big\{u(c)+\partial_a
v_j\,\big(y_j+r(t)a-c\big)\Big\}+\lambda_j\,(v_{-j}-v_j)
\qquad\text{on }(\uba,\infty)\times(0,T),
\end{equation}
with terminal condition $v_j(\cdot,T)=v_{T,j}$ and the
\emph{state-constraint} boundary condition
\begin{equation}\label{eq:BC}
\partial_a v_j(\uba,t)\ \ge\ u'\big(y_j+r(t)\uba\big),
\qquad j=1,2,
\end{equation}
which says that the marginal value of wealth at the borrowing limit is
at least the marginal utility of the consumption level
$\ubc_j(t):=y_j+r(t)\uba$ that just keeps the household at the limit;
equivalently, that chosen saving at the limit is nonnegative, so the
constraint is never violated. Optimal consumption satisfies
$c_j=(u')^{-1}(\partial_a v_j)$ in the interior and
$c_j(\uba,t)=\min\{(u')^{-1}(\partial_a v_j(\uba^+,t)),\,\ubc_j(t)\}$
at the limit; the saving policy is $s_j(a,t)=y_j+r(t)a-c_j(a,t)$.

The cross-sectional distribution of (income, wealth) is described by
measures $G_j(\cdot,t)$ on $[\uba,\infty)$ of total mass one, evolving
under the households' saving policies and the income switching by the
Kolmogorov forward equation, understood in a weak (measure) sense that
accommodates a point mass $m_1(t)$ of households \emph{at} the
borrowing limit (Definition~\ref{def:KFweak} in
Appendix~\ref{app:distribution}); $G_j(\cdot,0)=G_{0,j}$ is given.
This system---equations (12)--(17) of \citet{AHLLM2022}, together
with their market-clearing condition (4)---is the transition-dynamics
version of the continuous-time Huggett model
\citep{Huggett1993}.

\subsection{Equilibrium}
\label{sec:eqmdef}

\begin{definition}\label{def:eqm}
Given data $(G_0,v_T)$ and bond supply $B\in[0,\infty)$, an
\emph{equilibrium transition} on $[0,T]$ is a continuous interest-rate
path $r(\cdot)$, household value functions solving
\eqref{eq:HJB}--\eqref{eq:BC} with terminal data $v_T$ (in the
constrained viscosity sense, Definition~\ref{def:cvs}), and
distributions solving the forward equation from $G_0$ under the
induced policies, such that the bond market clears at every date:
\begin{equation}\label{eq:clearing}
S(t)\;:=\;\sum_{j=1,2}\int_{[\uba,\infty)}a\dd G_j(a,t)\;=\;B
\qquad\text{for all }t\in[0,T].
\end{equation}
\end{definition}

Throughout, candidate rate paths lie in a band
$\Kcal=\{r\in C([0,T]):r_{\min}\le r\le r_{\max}\}$ with
\begin{equation}\label{eq:band}
r_{\max}<\frac{y_1}{|\uba|},
\end{equation}
which guarantees $\ubc_1(t)=y_1+r(t)\uba\ge\ubc>0$: even the
low-income household at the borrowing limit can service its debt out
of income with something left to consume. Condition \eqref{eq:band} is
the transition version of the requirement that the borrowing limit be
tighter than the natural limit; it is what makes the constraint
economically meaningful (households at the limit are hand-to-mouth,
not defaulting) and, mathematically, it is the inward-pointing
condition on which the entire boundary analysis rests. The full list
of maintained assumptions on the data $(G_0,v_T)$ is collected in
Appendix~\ref{app:household}
(Assumptions~\ref{ass:prim}--\ref{ass:G0}, covering primitives, the
rate band, and the data);
the terminal valuation is concave, increasing, with marginal values
bounded above and decaying no faster than CRRA marginal utility of a
linearly growing consumption level---properties satisfied by the
canonical choice $v_\infty(\cdot\,;r^*)$.

For the quantitative equilibrium analysis we will distinguish two
regimes for the wealth distribution near the constraint:
\begin{itemize}[leftmargin=2.6em]
\item[\textbf{(D1)}] the absolutely continuous parts of the
distributions have densities bounded near $\uba$, uniformly in $t$
and over the band $\Kcal$;
\item[\textbf{(D1$'$)}] the densities may blow up like
$(a-\uba)^{-1/2}$ near $\uba$, with a constant uniform in $t$ and
over $\Kcal$.
\end{itemize}
(D1$'$) is the realistic regime: it is exactly the profile generated
by the square-root consumption law at a binding constraint
(Section~\ref{sec:binding}), just as in the stationary distribution of
\citet{AHLLM2022}. (D1) holds, for instance, transiently while the
constraint does not bind.

\subsection{Why transitions are not steady states strung together}
\label{sec:whyhard}

It is worth previewing, in words, where the difficulty lies, because
it is \emph{not} where the steady-state theory suggests.

In a stationary equilibrium the logic is: fix $r$, solve the
household problem, compute the implied stationary wealth stock
$S_{\mathrm{stat}}(r)$, and find $r$ with $S_{\mathrm{stat}}(r)=B$.
The map $r\mapsto S_{\mathrm{stat}}(r)$ is a well-behaved scalar
function, and existence follows from its boundary behavior
\citep{Aiyagari1994,AHLLM2022}. One might hope to treat a transition
as a family of such problems indexed by $t$. This fails, for an
economic reason: $S(t)$ is a \emph{stock}, accumulated from past
saving decisions, each of which was made looking at \emph{future}
rates. The rate at date $t$ influences wealth at date $t$ not at all;
it influences wealth \emph{after} $t$ through household budgets, and
consumption \emph{before} $t$ through anticipation. So the natural
analogue of Aiyagari's diagram---move $r(t)$ until $S(t)$ hits
$B$---moves nothing. Section~\ref{sec:clearing} makes this precise
and locates the two places where the contemporaneous rate does bite:
interest on the bond supply, and the budgets of constrained
households. Everything constructive in this paper is built on those
two levers; and the second lever is itself an equilibrium object,
switched on and off by the binding analysis of the next section.

%======================================================================
\section{Household behavior along an interest-rate path}
\label{sec:household}

This section takes the rate path $r\in\Kcal$ as given and establishes
what we need about individual behavior: that the household problem is
well posed at all (Theorem~\ref{thm:household}), and---the
economically substantive part---when households at the borrowing
limit stay there (Section~\ref{sec:binding}). Readers willing to take
well-posedness on faith can skim to
Section~\ref{sec:binding}; the technical package, including all
regularity and stability estimates used later, is developed in
Appendix~\ref{app:household}.

\begin{theorem}[the household problem is well posed]
\label{thm:household}
Under the maintained assumptions
(Appendix~\ref{app:household}), for every $r\in\Kcal$ the household's
value function is the unique constrained viscosity solution of
\eqref{eq:HJB}--\eqref{eq:BC} with terminal data $v_T$, in the class
of continuous functions concave and increasing in wealth with linear
growth. It is continuously differentiable in wealth on
$(\uba,\infty)$, with a one-sided derivative at $\uba$; consumption is
continuous and nondecreasing in wealth; marginal values admit
two-sided bounds, uniform over $\Kcal$,
\[
0<\kappa\,e^{-CT}(1+a-\uba)^{-\gamma}\ \le\ \partial_a v_j(a,t)\ \le\
\max\{\bar p_T,\,u'(\ubc)\}\,e^{(r_{\max}-\rho)^+T},
\]
where $\bar p_T$ bounds the terminal marginal value; equivalently,
consumption is bounded below by a positive constant and grows at most
linearly in wealth. The value function and the policies depend
continuously---and the value Lipschitz-continuously---on the rate
path, with only the \emph{future} of the path entering
(Lemma~\ref{lem:stability}).
\end{theorem}

Three comments on the economics of the proof (details in
Appendix~\ref{app:household}). First, the bounds on marginal value are
where finite horizon differs from the steady state. Stationarily, the
binding constraint supplies the boundary marginal value
\emph{exactly}---$\partial_a v_1(\uba)=u'(y_1+r\uba)$ whenever
$r<\rho$---and concavity does the rest. Along a transition, binding is
not guaranteed (next subsection), and the bounds must be built by
comparison with explicit sub- and super-solutions: the upper bound
propagates backward the larger of the terminal marginal value and the
boundary obstacle $u'(\ubc)$; the lower bound is a CRRA-shaped barrier
expressing that consumption cannot outgrow wealth linearly. The cap
$\bar p_T$ on terminal marginal values is a genuine assumption on the
terminal data, not a free good. Second, the bounds keep marginal
values in a compact positive range, which is what tames the CRRA
Hamiltonian $H(p)=\max_c\{u(c)-pc\}$, whose slope
$H'(p)=-p^{-1/\gamma}$ blows up as $p\downarrow0$ (marginal values
near zero correspond to unboundedly rich households). Third---an
honest simplification---neither $\rho>0$ nor $r<\rho$ appears
anywhere: on a finite horizon the terminal valuation replaces
transversality, and impatience plays no coercive role. The conditions
involving $\rho$ reappear only through the terminal data, which is
exactly where the next subsection locates them.

\subsection{When does the borrowing constraint bind?}
\label{sec:binding}

Say the constraint \emph{binds} at date $t$ if households at the
limit stay there: $s_1(\uba,t)=0$, equivalently
$\partial_a v_1(\uba^+,t)=u'(\ubc_1(t))$. When it binds, the point
mass $m_1(t)$ of hand-to-mouth households persists and grows; when it
fails ($s_1(\uba,t)>0$), the constrained group saves its way off the
limit and the mass disperses. In the stationary economy binding is
automatic under impatience, $r<\rho$. Along a transition the
governing comparison is different:

\begin{lemma}[binding near the end of the horizon]\label{lem:binding}
(i) If the terminal marginal value at the limit exceeds the marginal
utility of hand-to-mouth consumption at the terminal rate---
$\partial_a v_{T,1}(\uba)>u'(y_1+r(T)\uba)$---then the constraint does
\emph{not} bind on an interval $(T-\delta,T)$, regardless of how
$r(t)$ compares to $\rho$. (ii) If the reverse strict inequality
holds, the constraint binds near $T$.
\end{lemma}

\begin{corollary}[the knife edge at the canonical terminal data]
\label{cor:knife}
Let $v_T=v_\infty(\cdot\,;r^*)$, the stationary value at rate
$r^*$, so that $\partial_a v_{T,1}(\uba)=u'(y_1+r^*\uba)$ by the
stationary binding equality. Then the boundary saving rate at the end
of the horizon is
\[
s_1(\uba,T^-)\;=\;\big[(r(T)-r^*)\,\uba\big]^+,
\]
and since $\uba<0$: \emph{the constraint binds at $T^-$ if and only if
$r(T)\ge r^*$}. Whenever the path ends below the rate embedded in the
continuation values, the hand-to-mouth group disperses near the end of
the transition---even if $r(t)<\rho$ throughout.
\end{corollary}

The economics bears repeating because it inverts a steady-state
intuition. Households at the borrowing limit are \emph{debtors}: their
consumption is income net of debt service, $y_1+r(t)\uba$, which
\emph{falls} as the rate rises. Their continuation plan, embedded in
$v_T$, was drawn at rate $r^*$. If the transition delivers a terminal
rate below $r^*$, debt service is lighter than the continuation plan
assumed; the constrained household finds itself with slack, saves, and
detaches from the limit. Impatience ($\rho$) never enters: over a
short remaining horizon, behavior is governed by the terminal
valuation, not by discounting. A closed-form illustration with no
income risk, in which binding fails on an entire terminal interval
despite $r<\rho$, is given in Example~\ref{ex:closedform}
(Appendix~\ref{app:household}); the same example exhibits the turnpike
property---binding is restored at dates far from $T$, where the
stationary logic reasserts itself.

Near the constraint, consumption follows a square-root law, as in the
stationary analysis of \citet{AHLLM2022}, but with a time-dependent
speed:

\begin{lemma}[square-root law along a transition]\label{lem:sqrt}
Let $r\in C^1$ and suppose the constraint binds on
$[t_0,t_1]\subset(0,T)$. Define
\[
\nu_1(t)\;:=\;\frac{(\rho+\lambda_1-r(t))\,u'(\ubc_1(t))
-\lambda_1\,\partial_a v_2(\uba,t)}{-u''(\ubc_1(t))}
\;+\;\dot r(t)\,\uba .
\]
If $\nu_1\ge\ubar\nu>0$ on $[t_0,t_1]$, then locally uniformly there
\[
c_1(a,t)=\ubc_1(t)+\sqrt{2\nu_1(t)}\,\sqrt{a-\uba}+O(a-\uba),
\qquad
s_1(a,t)\sim-\sqrt{2\nu_1(t)}\,\sqrt{a-\uba}.
\]
\end{lemma}

Relative to the stationary coefficient, time dependence adds the term
$\dot r(t)\,\uba$: since $\uba<0$, a \emph{rising} rate path lowers
$\nu_1$---anticipated increases in debt service slow the approach to
the constraint---and $\nu_1(t)\downarrow0$ is precisely the incipient
dispersal regime of Corollary~\ref{cor:knife}. The formal expansion
and its rigorous version are discussed in
Appendix~\ref{app:household}.

Finally, the constrained mass itself. While the constraint binds,
$m_1(t)$ grows by the inflow of low-income households hitting the
limit and shrinks only as its members draw the high income state (at
rate $\lambda_1$). When binding fails, the entire mass detaches
\emph{at once} and moves into the interior as a traveling atom:
$m_1(t)$ drops to zero discontinuously
(Proposition~\ref{prop:atom}, Appendix~\ref{app:distribution}). This
bookkeeping matters beyond descriptive interest: $m_1(t)$ is exactly
the market's contemporaneous lever on the interest rate, so the
binding analysis of this subsection switches the tractability of the
equilibrium problem on and off.

%======================================================================
\section{Market clearing along a transition}
\label{sec:clearing}

We now turn to equilibrium. Throughout, $S(t)$ denotes aggregate
wealth \eqref{eq:clearing}, $C(t)=\sum_j\int c_j\dd G_{j,t}$ aggregate
consumption, and $Y(t)=\sum_j\int y_j\dd G_{j,t}$ aggregate income.

\subsection{Compatibility and the flow form of clearing}

\begin{proposition}[the data must clear the market at date zero]
\label{prop:compat}
If an equilibrium transition exists, then $S(0)=\sum_j\int a\dd
G_{0,j}=B$. No choice of prices can repair incompatible initial data:
clearing at $t=0$ involves the initial distribution alone.
\end{proposition}

This is obvious once stated, but it is a substantive discipline on
policy experiments: an ``MIT shock'' that revalues wealth or bonds at
$t=0$ must do so consistently with \eqref{eq:clearing}, or no
equilibrium path exists at all.

\begin{lemma}[clearing in stocks equals clearing at date zero plus
clearing in flows]\label{lem:flow}
Along any solution of the distributional dynamics with uniformly
bounded first moments, $t\mapsto S(t)$ is Lipschitz and
\[
\dot S(t)\;=\;\sum_j\int s_j\dd G_{j,t}\;=\;Y(t)+r(t)S(t)-C(t),
\]
the boundary contributing nothing (households at the limit have zero
saving while the constraint binds; dispersing mass is counted in the
interior). Consequently, for $r\in\Kcal$:
\[
S(t)=B\ \ \forall t\in[0,T]
\qquad\Longleftrightarrow\qquad
S(0)=B\ \text{ and }\ C(t)=Y(t)+r(t)B\ \ \forall t\in[0,T].
\]
\end{lemma}

The flow form is the economy's aggregate budget identity: consumption
absorbs income plus the interest paid on the bond supply. For $B=0$ it
says aggregate consumption is pinned to aggregate income at every
date---an exogenous restriction in which, as we will see, the interest
rate barely appears.

\subsection{Stock-level clearing is degenerate}

\begin{theorem}[no contemporaneous leverage in stocks]
\label{thm:firstkind}
Fix data and a path $r\in\Kcal$, and consider the equilibrium map
$r\mapsto S(\cdot\,;r)$.
\begin{enumerate}[label=(\roman*),leftmargin=2.2em]
\item \emph{Pinning.} $S(0;r)$ is independent of $r$: near $t=0$ the
stock condition restricts only the data.
\item \emph{Vanishing diagonal.} For path perturbations $\delta r$
supported on a window $(t_0-h,t_0+h)$,
\[
|S(t_0;r+\delta r)-S(t_0;r)|\ \le\ C\,h\,(1+T)\,\|\delta r\|_\infty
\]
(under the bounded-density regime (D1); with an additional factor
$1+\log_+\!\big(1/(h\|\delta r\|_\infty)\big)$ in the realistic
boundary-layer regime (D1$'$)---either way the bound vanishes with
$h$): the rate at date $t_0$ has no contemporaneous effect on wealth
at date $t_0$.
\item \emph{Compactness.} The linearized map $DS(\cdot\,;r)$ is a
compact integral operator on $C([0,T])$; the equation
$S(\cdot\,;r)=B$ is a first-kind equation, admitting no bounded local
inverse and no Picard or Newton theory, and a stock-level
t\^atonnement $r\mapsto r+\kappa(S-B)$ cannot be closed: unlike the
stationary problem, sign information on excess supply at the edges of
the rate band is unavailable pointwise in $t$, because $S(t;\cdot)$ is
a functional of the entire path.
\end{enumerate}
Any existence argument must therefore run through the flow form of
Lemma~\ref{lem:flow}.
\end{theorem}

The economic content of (ii) is the observation of the introduction:
consumption at $t$ is forward-looking, wealth at $t$ is
backward-looking, and the instant $t$ itself has measure zero in
both. The exception is hidden in the constant $C$: it is built from
the \emph{integrated} response of consumption and savings to the rate
window, which includes the mechanical response of constrained
households. That response reappears, no longer averaged away, in the
flow form: differentiating the flow-clearing residual with respect to
the contemporaneous rate,
\begin{equation}\label{eq:diagonal}
\frac{\partial}{\partial r(t)}\Big[C(t)-Y(t)-r(t)B\Big]
\;=\;m_1(t)\,\uba\;-\;B\;=\;-\big(B+|\uba|\,m_1(t)\big),
\end{equation}
because the only households whose date-$t$ consumption depends on the
date-$t$ rate are those at the borrowing limit, consuming
$y_1+r(t)\uba$. The flow constraint has a uniformly invertible
diagonal if and only if $B+|\uba|m_1(t)$ is bounded away from
zero---automatically for $B>0$; through the constrained mass, hence
through the binding analysis of Section~\ref{sec:binding}, when
$B=0$.

Theorem~\ref{thm:firstkind} also rationalizes numerical practice. The
transition algorithm of \citet[\S5.5]{AHLLM2022} updates
$r^{\ell+1}(t)=r^\ell(t)-\xi(t)\,\frac{d}{dt}S^\ell(t)$: the update
signal is the \emph{flow} imbalance, the object with a nondegenerate
diagonal, not the excess stock. Likewise, in discrete-time
sequence-space methods \citep{AuclertEtAl2021}, our result implies
that the Jacobian of stock-level clearing becomes ill-conditioned as
the period shrinks, while flow-level clearing does not.

%======================================================================
\section{Short transitions: existence and uniqueness}
\label{sec:smallT}

We now solve the equilibrium problem over short horizons, first with
positive bond supply, then at zero net supply. Throughout:
$S(0)=B$ (Proposition~\ref{prop:compat}), the maintained assumptions,
the household package of Section~\ref{sec:household}, and the two
density regimes (D1)/(D1$'$) of Section~\ref{sec:eqmdef}.

\subsection{Positive bond supply: the multiplier map and the
hand-to-mouth threshold}

For $B>0$, Lemma~\ref{lem:flow} converts equilibrium into a fixed
point over rate paths: $r=\Psi(r)$ with
\[
\Psi(r)(t)\;:=\;\frac{C(t;r)-Y(t;r)}{B},
\]
the rate that makes bond interest absorb the excess of consumption
over income implied by the path $r$. The map $\Psi$ is compact (its
range is uniformly bounded and equicontinuous---there is no identity
part, unlike a stock-level t\^atonnement), and its Lipschitz constant
decomposes along economic channels: for paths $r,r'$ with
$\eta=\|r-r'\|_\infty$,
\[
\|\Psi(r)-\Psi(r')\|_\infty\ \le\
\underbrace{\theta\,\eta}_{\substack{\text{contemporaneous:}\\
\text{hand-to-mouth}\\ \text{budgets}}}
\;+\;\underbrace{C_1T\,\eta}_{\substack{\text{anticipation:}\\
\text{policies respond}\\ \text{to future rates}}}
\;+\;\underbrace{C_2T\,\eta\,\big(1+\log_+(1/\eta)\big)}
_{\substack{\text{history:}\\ \text{distribution shaped}\\
\text{by past rates}}},
\]
where
\[
\theta\;:=\;\frac{|\uba|\,\sup_{t,\,r\in\Kcal}m_1(t;r)}{B},
\]
with the logarithm absent under (D1). The anticipation and history
channels are weak over short horizons---each carries a factor
$T$---but the contemporaneous channel is not: it is the hand-to-mouth
budget response, present at any horizon, and it is governed by
$\theta$, the ratio of constrained households' debt exposure to the
bond supply. One further data condition is needed to keep the fixed
point inside the rate band: as $T\to0$ the map $\Psi$ converges to an
affine map of the instantaneous rate (the affine term again being the
hand-to-mouth budget), whose fixed point $r^{**}$ solves a scalar
equation in the data; we require $r^{**}$ to be interior to the band
(condition (ANCH), Appendix~\ref{app:proofs}).

\begin{theorem}[existence and uniqueness for short transitions,
$B>0$]\label{thm:smallT}
Let $B>0$, $S(0)=B$, (ANCH), and
\[
\theta\ \le\ \bar\theta\ <\ 1
\qquad\text{\emph{(hand-to-mouth debt exposure below bond supply).}}
\]
Then there is $T^*>0$, depending on the data, $\bar\theta$, and the
anchoring margin, such that for $T\le T^*$:
\begin{enumerate}[label=(\roman*),leftmargin=2.2em]
\item an equilibrium transition exists;
\item under (D1), the map $\Psi$ is a contraction with modulus
$\bar\theta+CT<1$: the equilibrium is unique in the band, and the
natural iteration---guess a rate path, solve households backward,
aggregate forward, update the rate from the flow-clearing
condition---converges geometrically from any starting path;
\item under (D1$'$), any two equilibria are within
$\exp\!\big(-(1-\bar\theta-(C_1+C_2)T)/(C_2T)\big)$ of each other:
uniqueness holds up to an ambiguity exponentially small in $1/T$.
\end{enumerate}
\end{theorem}

The stability condition $\theta<1$ deserves emphasis; it is, to our
knowledge, new. Consider a contemporaneous increase $dr$ in the rate.
Flow clearing requires the economy to absorb $B\,dr$ of additional
interest income as consumption. But the same increase \emph{taxes}
hand-to-mouth debtors, whose consumption falls mechanically by
$|\uba|\,m_1\,dr$. If the second effect exceeds the first, the
market-clearing response to excess consumption is to \emph{raise}
required consumption further: the multiplier loop is unstable, and the
fixed-point argument fails. The condition is economically
interpretable and, in principle, measurable: the debt of
constrained households relative to the outstanding stock of liquid
assets. Economies with large hand-to-mouth populations, deep
borrowing, and thin asset markets are the ones in which transition
equilibria are fragile.

A natural question is whether the threshold is an artifact of the
fixed-point scheme. It is not, in the following sense: one can solve
the flow-clearing condition \emph{exactly} for $r(t)$ (the constrained
group's budget is known in closed form), eliminating the
contemporaneous channel altogether; but the resulting map depends on
the constrained mass $m_1(t;r)$ \emph{separated} from the rest of the
distribution, and the mass alone is only $\tfrac12$-H\"older stable in
the rate path under (D1$'$)---mass exchanged between the atom and the
adjacent boundary layer moves the aggregate only at second order, but
moves the atom at first order. Under (D1) the elimination works and
$\theta<1$ can be dispensed with; under the realistic (D1$'$) the
threshold is load-bearing (Proposition~\ref{prop:diagelim},
Appendix~\ref{app:proofs}).

\subsection{Zero net supply: the constrained mass is the market}
\label{sec:B0}

At $B=0$ the flow condition reads $C(t)=Y(t)$: aggregate consumption
is pinned to aggregate income, and by \eqref{eq:diagonal} the
\emph{only} contemporaneous lever of the rate on clearing is the
hand-to-mouth group, with weight $|\uba|\,m_1(t)$.

\begin{theorem}[$B=0$ with a constrained population]\label{thm:B0atom}
Let $B=0$, $S(0)=0$, $m_1(0)>0$, and assume---as an explicit
hypothesis---that the constraint binds along the band throughout
$[0,T]$. (Necessary for this is the terminal-data compatibility
$\partial_a v_{T,1}(\uba)\le u'(y_1+r_{\min}\uba)$, which for
canonical terminal data reads $r_{\min}\ge r^*$ by
Corollary~\ref{cor:knife}; but terminal compatibility controls binding
only near $T$, and interior dispersal episodes along
dipping-and-rebounding paths must be excluded, e.g.\ by restricting to
slowly varying paths on which the square-root coefficient $\nu_1$ of
Lemma~\ref{lem:sqrt} stays positive.) Assume also the $B=0$ anchoring
condition: the instantaneous rate implied by the data---the solution
of the flow-clearing equation at $t=0$, with policies at their
terminal-implied values---is interior to the band. Then the
constrained mass decays no faster than its members find high income,
$m_1(t)\ge m_1(0)e^{-\lambda_1t}$, the flow condition can be solved
for the rate, and for $T\le T^*(\text{data},m_1(0))$ an equilibrium
transition exists; it is unique in the band under (D1). The horizon
$T^*$ shrinks linearly with $m_1(0)$: a thin hand-to-mouth population
supports only a thin theorem.
\end{theorem}

\begin{proposition}[$B=0$ with no constrained population: no local
theory]\label{prop:notheorem}
Suppose $B=0$ and $m_1\equiv0$ on some interval---for instance near
$t=0$ when the initial distribution has no mass at the limit, or
following a dispersal episode. On that interval the clearing
condition $C(t;r)=Y(t)$ contains the contemporaneous rate
\emph{nowhere}: its linearization in the path is a compact operator
(as in Theorem~\ref{thm:firstkind}), the constraint is a first-kind
equation, and no local solution theory---Picard, Newton, or implicit
function---applies. In this regime relaxations are a necessity, not a
computational convenience: an $\varepsilon$-elastic bond supply
$B_\varepsilon(r)=\varepsilon r$, or a price rule
$r(t)=F(C(t)-Y(t))$ with $F$ Lipschitz, restores an invertible
diagonal $\varepsilon+|\uba|m_1$ and places the problem back in the
class of Theorem~\ref{thm:smallT}; as $\varepsilon\to0$ one retains
only weak limits of approximate equilibria, with clearing recovered in
an integrated sense. Existence of exact equilibria in this regime is
open---and the neighbouring analysis of \citet{Ambrose2021}, in a
model where the constrained-mass channel is absent by construction,
shows that there the exact problem is generically \emph{unsolvable}
while its relaxation is well posed over short horizons.
\end{proposition}

The pair of results delivers a sharp economic message about the
textbook case $B=0$: \emph{the hand-to-mouth population is the
market}. When it is present, its budget is the margin along which the
interest rate clears the market instant by instant, and equilibrium
transitions exist. When it is absent, the interest rate is a shadow
price with no contemporaneous market to clear, exact instant-by-instant
clearing overdetermines the system, and the economically meaningful
formulation is one in which either the asset supply has some
elasticity or the price responds to imbalances with some friction.

%======================================================================
\section{Long transitions: monotonicity and multiple self-fulfilling
paths}
\label{sec:long}

The results above are local in the horizon: the anticipation and
history channels each carry a factor $T$, and the ``loop gain'' of
expectations---beliefs about future rates alter consumption today,
which alters the wealth distribution, which alters market-clearing
rates tomorrow---scales like $T^2$. For long horizons, one hopes for
global structure. The natural candidate, suggested by the stationary
theory, is monotonicity in the rate.

\begin{lemma}[pathwise comparison fails for debtors]\label{lem:nocomp}
Couple two economies under rate paths $r_h\ge r_\ell$ (pointwise),
identical income realizations, identical initial wealth. Wealth under
the higher path dominates wealth under the lower path only if,
wherever wealth is negative,
\[
c^{r_\ell}(a,t)-c^{r_h}(a,t)\ \ge\ (r_h(t)-r_\ell(t))\,|a| :
\]
the consumption cut induced by higher rates must exceed the added
interest burden. For debtors this is a quantitative restriction on the
consumption response, not a sign restriction; ``higher rates
$\Rightarrow$ higher wealth'' is not a comparison-principle fact in
this model.
\end{lemma}

In the stationary economy, \citet{AHLLM2022} prove the aggregate
version of this dominance---saving increasing in $r$---when the
intertemporal elasticity of substitution is at least one
($\gamma\le1$) and borrowing is prohibited, and conclude uniqueness of
the steady state. Their argument does not survive transplantation to
time-varying paths, for identifiable reasons: it perturbs by a
\emph{constant} rate (a time-varying perturbation turns the linearized
Euler system into a two-point boundary problem with sign-indefinite
kernels mixing substitution and wealth effects); it uses CRRA
homogeneity to absorb the rate shift (which requires terminal data
covariant under the same scaling---and terminal data drawn from a
\emph{different} stationary equilibrium are not); and it aggregates
using stationary-distribution identities unavailable along a
transition. We therefore state the needed property as a hypothesis and
develop its consequences honestly as such:

\begin{itemize}[leftmargin=2em]
\item[\textbf{(H-MONO)}] For paths $r\le r'$ pointwise in the band:
$C(t;r)\ge C(t;r')$ for all $t$ (aggregate consumption falls when the
whole rate path rises).
\end{itemize}

\begin{theorem}[monotone localization of equilibria]\label{thm:mono}
Let $B>0$, $S(0)=B$, assume the multiplier map preserves the band and
(H-MONO). Then $\Psi$ is antitone, and the two-sided iteration started
from the band edges,
$\bar r^{n+1}=\Psi(\ubar r^{n})$, $\ubar r^{n+1}=\Psi(\bar r^{n})$,
converges monotonically and uniformly to continuous limits
$\ubar r^\ast\le\bar r^\ast$ with $\Psi(\ubar r^\ast)=\bar r^\ast$
and $\Psi(\bar r^\ast)=\ubar r^\ast$. \emph{Every} equilibrium in the
band lies between $\ubar r^\ast$ and $\bar r^\ast$; any two
pointwise-ordered equilibria coincide; and if the limits collapse,
$\ubar r^\ast=\bar r^\ast$, the equilibrium is unique and the
iteration converges to it globally.
\end{theorem}

The bracket $[\ubar r^\ast,\bar r^\ast]$ is computable: it is the
limit of the natural best-response iteration started from the
extremes, and the collapse condition can be checked numerically. For
$\gamma\le1$, (H-MONO) is credible---its stationary counterpart holds
under a no-borrowing limit
by \citet{AHLLM2022}, and over short horizons up to $O(T)$ by the
sensitivity decomposition---but we emphasize it remains unproven for
genuinely time-varying perturbations.

For $\gamma>1$ the situation is qualitatively different.

\begin{proposition}[risk aversion above one: singular linearizations
and multiplicity]\label{prop:multiplicity}
By \citet[Thm.~1.10]{Hofer2025}, for every $\gamma>1$ there exist
two-state economies (with a no-borrowing limit) and rates
$r_1<r_2<\rho$ at which stationary aggregate wealth satisfies
$S_{\mathrm{stat}}(r_1)>S_{\mathrm{stat}}(r_2)$. Two scope caveats:
the construction is existential over economies (a rare-disaster
income process), not universal; and it is proved at the no-borrowing
limit $\uba=0$, a boundary case of our maintained $\uba<0$ (where the
hand-to-mouth channel $m_1\uba$ vanishes identically), so its
extension to $\uba<0$ rests on an unproven continuity argument. On
the constructed class, and with the bridge from constant paths to
transitions understood as a turnpike heuristic:
\begin{enumerate}[label=(\roman*),leftmargin=2.2em]
\item (H-MONO) fails in the bulk of long horizons: constant paths at
$r_1<r_2$ order aggregate wealth the wrong way for most of the
horizon.
\item At a stationary point on the downward-sloping branch, slowly
varying path perturbations satisfy $DS[\delta r]\approx
S_{\mathrm{stat}}'(r^\ast)\,\delta r$ with
$S_{\mathrm{stat}}'(r^\ast)\le0$: the linearized clearing map acquires
an approximate kernel (at a fold) or a wrong-signed diagonal. Monotone
methods and na\"ive continuation in the horizon cross a singular
linearization; the reduction at a fold is one-dimensional and
degenerate---the classical environment for bifurcation of
self-fulfilling equilibrium paths, in which a believed future path of
rates alters current consumption, reshapes the wealth distribution,
and validates the belief.
\item When several stationary rates clear the same bond supply, data
connecting two of them (initial distribution from one, terminal
values from the other; both satisfy $S(0)=B$) admit, on long
horizons, transition paths lingering near each steady state as natural
candidates for \emph{distinct equilibria with identical endpoints}.
Under $\gamma\le1$ and (H-MONO), by contrast,
Theorem~\ref{thm:mono} confines any multiplicity to non-ordered,
crossing pairs of paths.
\end{enumerate}
\end{proposition}

\begin{remark}[what is proved and what is conjectured]
\label{rem:honest}
Unconditional, given the household package:
Proposition~\ref{prop:compat}, Lemma~\ref{lem:flow},
Theorem~\ref{thm:firstkind}, Theorems~\ref{thm:smallT} and
\ref{thm:B0atom}, Proposition~\ref{prop:notheorem}. Conditional:
Theorem~\ref{thm:mono} under (H-MONO)---credible for $\gamma\le1$,
refuted in the bulk-of-horizon sense for $\gamma>1$ on the economies
of \citet[Thm.~1.10]{Hofer2025}. Open: existence for long horizons at
$B>0$ without (H-MONO) (the missing ingredient is an a~priori bound
keeping equilibrium rate paths in a band); exact uniqueness in the
boundary-layer density regime; existence at $B=0$ without a
constrained population; and the rigorous turnpike bridge in
Proposition~\ref{prop:multiplicity}.
\end{remark}

%======================================================================
\section{Conclusion}
\label{sec:conclusion}

The steady-state theory of heterogeneous-agent economies rests on two
gifts: the interest rate summarizes the distribution in one number,
and impatience pins households to the borrowing constraint, where
their behavior is transparent. This paper shows that along a
transition both gifts are revoked, and identifies exactly what
replaces them. Market clearing in stocks is degenerate---today's rate
moves nothing today---and the operative equilibrium condition is the
aggregate budget identity in flows, in which the rate's only
instantaneous levers are interest on the bond supply and the budgets
of hand-to-mouth debtors. Binding at the constraint is governed by
terminal conditions rather than impatience, so the hand-to-mouth
population---the market's own lever---can disperse endogenously along
the path.

Three practical conclusions follow. First, transition equilibria
exist and are unique over short horizons under a measurable stability
condition, $\theta=|\uba|\sup_tm_1/B<1$: hand-to-mouth debt
below the stock of liquid assets. Second, the textbook case of zero
net supply is exactly as well behaved as its constrained population:
with no one at the limit, instant-by-instant clearing overdetermines
the economy, and slightly elastic supply or a price rule is not an
approximation of the true problem but its correct formulation---a
conclusion reinforced from the frictionless side by
\citet{Ambrose2021}. Third, for risk aversion above one, the
non-monotone aggregate-saving economies of \citet{Hofer2025} can make
the linearized clearing condition singular along a transition, and
multiple self-fulfilling paths between the same endpoints become a
live possibility: the uniqueness of computed transition dynamics in
common calibrations is an article of numerical faith, and testing for
it---by restarting path-iteration algorithms from dispersed initial
guesses, or by monitoring the sign of the flow-clearing
diagonal---seems a worthwhile discipline.

%======================================================================
\appendix

\section{The household problem: technical package}
\label{app:household}

This appendix collects the assumptions and the analytical results
behind Theorem~\ref{thm:household} and
Section~\ref{sec:binding}. Throughout, the Hamiltonian is
$H(p)=\sup_{c>0}\{u(c)-pc\}=\frac{\gamma}{1-\gamma}p^{1-1/\gamma}$
($\gamma\ne1$; $H(p)=-1-\log p$ for $\gamma=1$): finite only for
$p>0$ when $\gamma\le1$ (for $\gamma>1$, $H(0)=0$ but the maximizer
escapes), strictly convex, with $H'(p)=-p^{-1/\gamma}$ singular as
$p\downarrow0$ for every $\gamma$.

\subsection{Maintained assumptions}

\begin{assumption}[primitives]\label{ass:prim}
$\gamma>0$; $\rho\ge0$; $0<y_1<y_2$; $\lambda_1,\lambda_2>0$;
$-\infty<\uba<0$; $T<\infty$; $B\in[0,\infty)$.
\end{assumption}

\begin{assumption}[price band]\label{ass:band}
$r\in\Kcal=\{r\in C([0,T]):r_{\min}\le r\le r_{\max}\}$ with
$r_{\max}<y_1/|\uba|$; write $\bar r=\max\{|r_{\min}|,r_{\max}\}$ and
$\ubc=y_1+r_{\max}\uba>0$.
\end{assumption}

\begin{assumption}[terminal data]\label{ass:VT}
$v_{T,j}\in C^1([\uba,\infty))$, concave, strictly increasing, with
constants $c_T^{\min}>0$, $C_T>0$ such that
\emph{(VT-u)} $\partial_a v_{T,j}\le u'(c_T^{\min})=:\bar p_T$ and
\emph{(VT-l)} $\partial_a v_{T,j}(a)\ge u'(C_T(1+a-\uba))$ for all
$a$. Both hold for $v_T=v_\infty(\cdot\,;r^*)$ with
$r^*\in[r_{\min},r_{\max}]$, $r^*<\rho$, by the stationary binding
equality and asymptotically linear stationary consumption.
\end{assumption}

\begin{assumption}[initial data]\label{ass:G0}
$G_0=(G_{0,1},G_{0,2})$ nonnegative measures on $[\uba,\infty)$ with
total mass one and finite first moment; compact support
$[\uba,A_0]$ where stated. Density regimes (D1)/(D1$'$) as in
Section~\ref{sec:eqmdef}.
\end{assumption}

\begin{remark}[not assumed]
Neither $\rho>0$, nor $r(t)<\rho$, nor the stationary finiteness
condition $\rho>(1-\gamma)r$ appears in this appendix. On a finite
horizon the terminal valuation replaces transversality.
\end{remark}

\subsection{Value function and constrained viscosity solutions}

The value function is
$v_j(a,t)=\sup\E[\int_t^Te^{-\rho(s-t)}u(c_s)ds
+e^{-\rho(T-t)}v_{T,y_T}(a_T)]$ over progressively measurable
$c\ge0$ maintaining $a_s\ge\uba$; admissibility is nonempty uniformly
over $\Kcal$ by Assumption~\ref{ass:band} (hold the state consuming
$\ubc_{y_s}(s)>0$). Standard arguments give pathwise wealth and
budget bounds, finiteness and linear growth of $v$, concavity and
strict monotonicity in $a$, and dynamic programming.

\begin{definition}[constrained viscosity solution]\label{def:cvs}
$(v_1,v_2)$, continuous on $[\uba,\infty)\times[0,T]$, is a
constrained viscosity solution if it is a viscosity subsolution of
\eqref{eq:HJB} on $(\uba,\infty)\times[0,T)$, a supersolution on
$[\uba,\infty)\times[0,T)$ (boundary included), and matches the
terminal data. For concave solutions the supersolution property at
$\uba$ is equivalent to \eqref{eq:BC} \citep{Soner1986a,CDL1990}.
\end{definition}

\begin{theorem}[existence, comparison, uniqueness]\label{thm:cvs}
Under Assumptions~\ref{ass:prim}--\ref{ass:VT}, for every
$r\in\Kcal$ the value function is the unique constrained viscosity
solution in the class of continuous functions concave, nondecreasing,
of linear growth in $a$.
\end{theorem}

\emph{Proof architecture.} (i) Sub/supersolution properties from
dynamic programming; the state-constraint supersolution property at
$\uba$ via Soner's inward cone, whose hypothesis is exactly
$\ubc_j(t)\ge\ubc>0$, uniform over $\Kcal$
\citep{Soner1986a,Soner1986b,CDL1990}. (ii) The singular Hamiltonian
is handled by truncation: comparison is proved for a globally
Lipschitz $H_\flat$ agreeing with $H$ on the compact gradient range
delivered by Theorems~\ref{thm:upper}--\ref{thm:lower}; solutions
never see the truncation. This interface between the gradient bounds
and the comparison argument must be made explicit, since the range is
not free in finite horizon. (iii) Comparison for the weakly coupled
system: the switching coupling is quasi-monotone, so the system
comparison principles of \citet{IshiiKoike1991} and
\citet{EnglerLenhart1991} combine with parabolic doubling
\citep{CIL1992} and a linear penalization at infinity; the boundary is
handled on the supersolution side by the inward cone, uniformly in
$t$. (iv) Verification along Filippov solutions of the closed-loop
dynamics (Appendix~\ref{app:distribution}). (v) $C^1$ regularity in
$a$: a concave viscosity solution cannot kink where the Hamiltonian
is strictly convex \citep{CannarsaSinestrari2004}; switching terms are
zeroth order and harmless.\hfill$\square$

\begin{remark}[terminal boundary layer]\label{rem:BClayer}
The boundary inequality \eqref{eq:BC} holds at every $t<T$ even if
the terminal data violate it at $t=T$: a marginal unit of wealth at
the limit can always be consumed immediately at rate
$\ge\ubc_j(t)$. If $\partial_a v_{T,j}(\uba)<u'(\ubc_j(T))$, the
boundary gradient jumps as $t\uparrow T$---a terminal boundary layer
in $\partial_a v$, invisible in $v$.
\end{remark}

\subsection{Gradient bounds by barriers}

Formally, $w_j=\partial_a v_j$ solves the cooperative transport
system
\begin{equation}\label{eq:W}
\partial_t w_j+s_j\,\partial_a w_j
=(\rho+\lambda_j-r(t))\,w_j-\lambda_j w_{-j},\qquad
w_j(\cdot,T)=v'_{T,j};
\end{equation}
the rigorous versions of the following run the barriers on
vanishing-viscosity or implicit finite-difference approximations and
pass to the limit.

\begin{theorem}[upper bound]\label{thm:upper}
For all $j,a,t$ and $r\in\Kcal$:
$\partial_a v_j(a,t)\le p_{\max}
:=\max\{\bar p_T,u'(\ubc)\}\,e^{(r_{\max}-\rho)^+T}$.
\end{theorem}

\emph{Sketch.} $W(t)=\max\{\bar p_T,u'(\ubc)\}
\exp(\int_t^T(r-\rho)^+)$ is a spatially constant supersolution of
\eqref{eq:W} (the required inequality $-(r-\rho)^+W\le(\rho-r)W$
always holds); at the boundary the state-constraint mechanism can
raise $w_j(\uba,t)$ at most to the obstacle
$u'(\ubc_j(t))\le u'(\ubc)\le W(t)$; at $T$, $W\ge\bar p_T$.
Cooperative comparison concludes. What replaced the stationary
shortcut: stationarily $w_1(\uba)=u'(y_1+r\uba)$ exactly (binding
under $r<\rho$); here binding may fail
(Lemma~\ref{lem:binding}), the bound comes from the obstacle or the
terminal data, and the cap (VT-u) is a genuine
assumption.\hfill$\square$

\begin{theorem}[lower bound; linear consumption growth]
\label{thm:lower}
There are $\kappa>0$ and $C=C(\gamma,y_2,\bar r,\rho,\lambda)$ with
$\partial_a v_j(a,t)\ge\kappa e^{-CT}(1+a-\uba)^{-\gamma}$,
equivalently $c_j(a,t)\le\kappa^{-1/\gamma}e^{CT/\gamma}(1+a-\uba)$,
uniformly over $\Kcal$.
\end{theorem}

\emph{Sketch.} The ansatz $m(a,t)=\delta(t)(1+a-\uba)^{-\gamma}$,
identical across income states so the switching coupling cancels, is
a subsolution of \eqref{eq:W} for $\delta(t)=\kappa e^{-C(T-t)}$,
using the one-sided drift bound $s_j\le y_2+\bar
r|\uba|+r^+_{\max}(1+a-\uba)$; at $T$ it sits below $v'_{T,j}$ by
(VT-l); at the boundary below the obstacle.\hfill$\square$

\subsection{Binding, dispersal, and the square-root law}

Lemma~\ref{lem:binding} follows from the stability of $v$ near $T$
plus concavity: as $t\uparrow T$,
$\partial_a v_1(\uba^+,t)\to\max\{\partial_a
v_{T,1}(\uba),u'(\ubc_1(T))\}$, and comparing the implied desired
consumption with hand-to-mouth income $\ubc_1(t)$ yields the two
cases; Corollary~\ref{cor:knife} is the specialization to canonical
terminal data.

\begin{example}[closed form: binding fails although $r<\rho$]
\label{ex:closedform}
Remove income risk ($\lambda_1=\lambda_2=0$, income $y$), let
$r<\rho$ be constant, and take terminal data
$v_T(a)=\omega_T^{\gamma}u(a+h_T)$ with $h_T>-\uba$. The problem
solves in closed form: $c(a,t)=(a+h(t))/\omega(t)$ with
$h(t)=\int_t^Ty\,e^{-r(s-t)}ds+h_Te^{-r(T-t)}$ (human wealth) and
$\dot\omega=\frac{\rho-(1-\gamma)r}{\gamma}\omega-1$,
$\omega(T)=\omega_T$ (the inverse marginal propensity to consume).
The boundary saving rate $s(\uba,t)=y+r\uba-(\uba+h(t))/\omega(t)$
is strictly positive at $t=T$ for $\omega_T$ large (a patient,
high-valuation terminal datum): binding fails on a whole terminal
interval despite $r<\rho$. As $T-t\to\infty$, $h$ and $\omega$
approach their stationary limits and binding is restored: the
finite-horizon boundary behavior agrees with the stationary one only
at distance $O(1)$ from $T$ (a turnpike statement).
\end{example}

\emph{On Lemma~\ref{lem:sqrt}.} Matching orders in \eqref{eq:W} at
the boundary along a binding interval, with
$w_1(\uba,t)=u'(\ubc_1(t))$ so that
$\partial_tw_1(\uba,t)=u''(\ubc_1(t))\,\dot r(t)\uba$, yields the
displayed $\nu_1(t)$; the singular products cancel exactly as in the
stationary case. The sign $\nu_1(t)\ge0$ coincides with the
complementarity (KKT) rate of the state-constraint multiplier, so the
square-root law and binding degrade together. A rigorous version
freezes coefficients on short time intervals and compares with sub-
and super-boundary layers; we flag it as the delicate step of the
package.

\subsection{Stability in the rate path}

\begin{lemma}[continuity and Lipschitz stability]\label{lem:stability}
If $r_n\to r$ uniformly in $\Kcal$, then $v^{r_n}\to v^r$ locally
uniformly (half-relaxed limits plus comparison, the boundary obstacle
converging uniformly by Assumption~\ref{ass:band}), and policies
converge locally uniformly on $[\uba,\infty)\times[0,T)$. Moreover
\[
|v^r_j(a,t)-v^{r'}_j(a,t)|\ \le\
C^\star(1+a-\uba)\,\|r-r'\|_{C([t,T])},
\]
with $C^\star$ uniform over $\Kcal$; only the future of the path
enters.
\end{lemma}

\begin{remark}[policies across the boundary layer]
\label{rem:halfHolder}
Pointwise Lipschitz stability of policies in the rate fails near the
limit: the square-root layer's coefficient and its binding/release
times move with $r$, making $\|c^r-c^{r'}\|_\infty\lesssim
\|r-r'\|^{1/2}$ the honest generic rate. The equilibrium analysis
needs less: in aggregates, the atom's consumption
$y_1+r(t)\uba$ is exactly Lipschitz with constant $|\uba|$, and the
interior part meets the layer only on a set of small mass
($O(\delta^{3/2})$ under (D1), $O(\delta)$ under (D1$'$)); the
pointwise loss integrates away, up to a logarithm under (D1$'$).
This atom/interior division is the interface used in
Section~\ref{sec:smallT}.
\end{remark}

\section{Distributional dynamics}
\label{app:distribution}

\begin{definition}[measure-valued forward solutions]\label{def:KFweak}
Given policies $c_j$ and drifts $s_j$, a narrowly continuous family
$t\mapsto(G_{1,t},G_{2,t})$ of nonnegative measures on
$[\uba,\infty)$ solves the forward equation if for all
$\varphi_j\in C^1_b([\uba,\infty)\times[0,T])$---including test
functions \emph{not} vanishing at $\uba$---
\[
\frac{d}{dt}\sum_j\int\varphi_j\dd G_{j,t}
=\sum_j\int\Big[\partial_t\varphi_j+s_j\partial_a\varphi_j
+\lambda_j(\varphi_{-j}-\varphi_j)\Big]\dd G_{j,t},
\qquad G_{j,0}=G_{0,j}.
\]
This is the time-dependent version of the extended weak formulation of
\citet{AHLLM2022} (Supplementary Appendix~E), written out there for
the two-type equation in stationary form with the time-dependent case
declared analogous; unbounded test functions such as $\varphi=a$ are
admitted by cutoff for solutions with uniformly bounded first
moments.
\end{definition}

\begin{lemma}[one-sided Lipschitz drift]\label{lem:OSL}
Since $c_j(\cdot,t)$ is nondecreasing,
$(s_j(a,t)-s_j(b,t))(a-b)\le\bar r|a-b|^2$: the drift is one-sided
Lipschitz uniformly in $t$ and $\Kcal$, despite the square-root
singularity (which has the contractive sign); and
$s_1(\uba,t)\ge0$ always.
\end{lemma}

\begin{theorem}[well-posedness of the forward equation]\label{thm:KF}
There is a unique measure-valued solution: the law of the
piecewise-deterministic process alternating the forward Filippov flows
of $\dot a=s_j(a,t)$ with the Poisson switching clock. Trajectories
merge at $\uba$ (atom formation) but never cross; uniqueness follows
by duality in the one-dimensional one-sided-Lipschitz transport
theory \citep{Filippov1960,BouchutJames1998,PoupaudRascle1997},
extended to the switching system by Duhamel. Compact supports
propagate with explicit bounds, and the first moment $S(t)$ is
Lipschitz in $t$ uniformly over $\Kcal$---the rigorous content of the
flow identity in Lemma~\ref{lem:flow}.
\end{theorem}

\begin{proposition}[atom dynamics: growth and discontinuous release]
\label{prop:atom}
Let $m_1(t)=G_{1,t}(\{\uba\})$. On binding intervals $m_1$ is
absolutely continuous with
\[
\dot m_1(t)=\Phi_1(t)-\lambda_1m_1(t),\qquad
\Phi_1(t)=-\lim_{a\downarrow\uba}s_1(a,t)g_1(a,t)\ \ge0,
\]
inflow through the square-root drift feeding the mass, leakage only
through switching to high income. At a release time $t_0$ (binding
fails just after $t_0$, e.g.\ the rate crosses below the
terminal-implicit rate of Corollary~\ref{cor:knife}), the mass
detaches \emph{as a whole} and travels into the interior along the
Filippov characteristic through $(\uba,t_0)$, decaying at rate
$\lambda_1$; hence $m_1$ drops to zero discontinuously, and
re-accumulates if binding resumes. In particular $t\mapsto m_1(t)$ is
of bounded variation but not continuous, and the contemporaneous
clearing sensitivity $B+|\uba|m_1(t)$ of
Section~\ref{sec:clearing} is switched on and off by the binding
analysis of Section~\ref{sec:binding}.
\end{proposition}

\section{Proof architectures for
Sections~\ref{sec:clearing}--\ref{sec:long}}
\label{app:proofs}

\paragraph{Lemma~\ref{lem:flow}.} Take $\varphi_j=a$ in
Definition~\ref{def:KFweak} (justified by cutoff): the switching
terms cancel, leaving $\dot S=\sum_j\int s_j\dd G_{j,t}
=Y+rS-C$. Given $S(0)=B$ and $C=Y+rB$, $\dot S=r(S-B)$, so
$S\equiv B$ by Gr\"onwall; conversely $S\equiv B$ forces $\dot S=0$,
i.e.\ $C=Y+rB$.

\paragraph{Theorem~\ref{thm:firstkind}.} (i) is
Proposition~\ref{prop:compat}. (ii): the value at dates $t\le t_0-h$
responds to a perturbation supported on $(t_0-h,t_0+h)$ only through a
Duhamel source active on a window of width $2h$
(Lemma~\ref{lem:stability}); the distribution responds directly only
during the window (drift term $a\,\delta r$); integrating the policy
response against the distribution over $[0,t_0]$ gives the bound,
with the logarithmic factor under (D1$'$) inherited from
Remark~\ref{rem:halfHolder}. (iii): (ii) plus Arzel\`a--Ascoli; the
first-kind character is the vanishing diagonal.

\subsection{Full proof of Theorem~\ref{thm:smallT}}
\label{app:smallTproof}

Throughout this subsection $B>0$, $S(0)=B$, the maintained
assumptions hold, distributions are supported in the common compact
$[\uba,A]$ with $A=A(T,A_0)$ from Theorem~\ref{thm:KF}, and for
$r,r'\in\Kcal$ we write $\eta:=\|r-r'\|_\infty$,
$\bar m_1:=\sup_{t,\,r\in\Kcal}m_1(t;r)$, and
$\theta=|\uba|\bar m_1/B\le\bar\theta<1$. The proof consumes the
household package of
Appendices~\ref{app:household}--\ref{app:distribution} through the
following quantitative statements, recorded once.

\begin{enumerate}[label=(S\arabic*),leftmargin=3em]
\item \emph{Policy stability.} At the constraint, whenever both
problems bind, the exact budget identity
$c^r_1(\uba,t)-c^{r'}_1(\uba,t)=\uba\,(r(t)-r'(t))$ holds (both sides
consume $y_1+r(\cdot)\uba$); in the mixed binding case the
$\min$-structure of the boundary policy gives
$|c^r_1(\uba,t)-c^{r'}_1(\uba,t)|\le|\uba|\eta$ plus a term of
interior type. On the interior region, uniformly over $\Kcal$,
\[
\sup_{a\in[\uba+\delta_0,A]}|c^r_j(a,t)-c^{r'}_j(a,t)|
\ \le\ L_1\,(T-t)\,\eta,
\]
obtained by differentiating the comparison argument behind
Lemma~\ref{lem:stability}---whose error term accrues only over
$[t,T]$, giving the factor $(T-t)$---on the region where the value is
locally $C^{1,1}$ in wealth. Across the boundary layer
$[\uba,\uba+\delta_0]$ the pointwise rate degrades to order
$\sqrt\eta$ (Remark~\ref{rem:halfHolder}); the Lipschitz rate is
restored there under uniform binding, where the square-root
coefficient $\nu_1(\cdot\,;r)$ of Lemma~\ref{lem:sqrt} is Lipschitz
in the path through its explicit formula. This layer statement is the
delicate step of the household package flagged in
Appendix~\ref{app:household}; it is used below only through
integrated quantities.
\item \emph{Layer modulus.} Uniformly over $\Kcal$,
$|c^r_j(a,t)-c^r_j(\uba,t)|\le C_{\mathrm{lay}}\sqrt{a-\uba}$ on
$[\uba,\uba+\delta_0]$ (Lemma~\ref{lem:sqrt} with the gradient bounds
of Theorems~\ref{thm:upper}--\ref{thm:lower}).
\item \emph{Flow stability.} $W_\infty(G^r_t,G^{r'}_t)\le
L_2\,t\,\eta$. \emph{Proof:} couple the two populations pathwise
(same initial wealth, same income realizations). Between switches,
$\tfrac{d}{dt}(a^r_t-a^{r'}_t)^2\le
2\bar r\,(a^r_t-a^{r'}_t)^2
+2|a^r_t-a^{r'}_t|\,\big(A\,\eta+\|c^r-c^{r'}\|_\infty\big)$ by
Lemma~\ref{lem:OSL} (the one-sided part has the contractive sign),
and the constrained projection at $\uba$ is a contraction in one
dimension; Gr\"onwall and (S1) give
$|a^r_t-a^{r'}_t|\le e^{\bar rt}\,t\,(A+|\uba|+L_1T)\,\eta$, and the
supremum over paths bounds $W_\infty$.
\item \emph{Time regularity.} $t\mapsto G_t$ is $W_\infty$-Lipschitz
with constant $\bar s=\sup_{[\uba,A]}|s_j|$ (bounded drift,
Theorem~\ref{thm:KF}), and away from the terminal layer the policies
have a uniform modulus of continuity in $t$ (bound $\partial_tv$
through the equation and the gradient bounds, then argue as in (S2)).
\end{enumerate}

\begin{lemma}[aggregate stability transfer]\label{lem:agg}
For all $t\in[0,T]$ and $r,r'\in\Kcal$:
\[
|C(t;r)-C(t;r')|\ \le\
|\uba|\,\bar m_1\,\eta+C_1T\,\eta+C_2T\,\eta\big(1+\log_+(1/\eta)\big)
\quad\text{under (D1$'$)},
\]
and $|C(t;r)-C(t;r')|\le(|\uba|\bar m_1+C_3T)\,\eta$ under (D1).
Moreover $Y(t;r)=Y(t)$ does not depend on the path: the income-type
marginals solve the autonomous two-state master equation
($\dot\pi_1=-\lambda_1\pi_1+\lambda_2\pi_2$), as is seen by taking
$\varphi_1\equiv1$, $\varphi_2\equiv0$ in
Definition~\ref{def:KFweak}.
\end{lemma}

\begin{proof}
Decompose
\[
C(t;r)-C(t;r')
=\underbrace{\sum_j\int\big(c^r_j-c^{r'}_j\big)(\cdot,t)\dd
G^{r'}_{j,t}}_{=:I}
+\underbrace{\sum_j\int c^r_j(\cdot,t)\dd\big(G^r_{j,t}
-G^{r'}_{j,t}\big)}_{=:II}.
\]

\emph{Term $I$.} The atom at $\uba$ contributes
$m_1(t;r')\,|c^r_1(\uba,t)-c^{r'}_1(\uba,t)|\le\bar m_1|\uba|\,\eta$
by the exact identity in (S1) (mixed binding case: same bound plus an
interior-type term absorbed into the next one). The absolutely
continuous part contributes, by (S1),
$\le L_1(T-t)\eta\cdot G^{r'}_t\big((\uba,A]\big)\le L_1T\eta$; the
boundary-layer portion of the integral is covered by the same
constant under the layer-restored Lipschitz rate of (S1), and in any
case by the integrated bound below.

\emph{Term $II$.} Let $\pi_t$ be an optimal $W_\infty$-coupling of
$G^{r}_{j,t}$ and $G^{r'}_{j,t}$ (in one dimension, the monotone
quantile coupling), so that pairs $(x,y)\in\operatorname{supp}\pi_t$
satisfy $|x-y|\le\hat\eta:=L_2T\eta$ by (S3). Then
$|II|\le\sum_j\int|c^r_j(x,t)-c^r_j(y,t)|\dd\pi_t(x,y)$. Split the
pairs three ways. (a) Both endpoints at the atom: zero contribution.
(b) Mismatch pairs (one endpoint at $\uba$, the other in the layer):
their $x$-coordinates lie within $\hat\eta$ of $\uba$, so by (S2) the
consumption discrepancy is at most
$C_{\mathrm{lay}}\sqrt{\hat\eta}$; under (D1$'$) the mass of the
layer $(\uba,\uba+\hat\eta]$ is at most
$D\int_0^{\hat\eta}x^{-1/2}dx=2D\sqrt{\hat\eta}$, so the mismatch
set contributes at most
$2DC_{\mathrm{lay}}\,\hat\eta=O(T\eta)$---mass exchanged between the
atom and the adjacent layer moves consumption only at second order,
because consumption is continuous across the constraint. (c) Interior
pairs: where both endpoints exceed $\uba+\hat\eta$, use the modulus
$|c^r_j(x)-c^r_j(y)|\le
C_{\mathrm{lay}}\,\hat\eta/\sqrt{x-\uba}$ for $x$ in the layer (mean
value of the $\sqrt{\cdot}$ profile) and plain Lipschitz continuity
of $c^r_j$ outside it; integrating against the (D1$'$) density,
\[
\int_{\hat\eta}^{\delta_0}\frac{\hat\eta}{\sqrt{x}}\;D\,
x^{-1/2}\dd x
\;=\;D\,\hat\eta\,\log\!\big(\delta_0/\hat\eta\big)
\;=\;O\!\big(T\eta\,(1+\log_+(1/\eta))\big),
\]
while the region outside the layer contributes
$\le\mathrm{Lip}(c)\,\hat\eta=O(T\eta)$. Summing (a)--(c) gives
$|II|\le C_2T\eta(1+\log_+(1/\eta))$ under (D1$'$); under (D1) the
density is bounded, every step above is Lipschitz, and
$|II|\le C\hat\eta=O(T\eta)$. Collecting terms yields the two
displayed bounds.
\end{proof}

\begin{lemma}[compact range]\label{lem:cpt}
$\Psi$ maps $\Kcal$ into a bounded, uniformly equicontinuous subset
of $C([0,T])$; its closed convex hull is compact.
\end{lemma}

\begin{proof}
Boundedness: $0\le C(t;r)\le\kappa^{-1/\gamma}e^{CT/\gamma}(1+A-\uba)$
by the linear consumption bound (Theorem~\ref{thm:lower}) on the
common support, so $\sup_{\Kcal}\|\Psi(r)\|_\infty<\infty$.
Equicontinuity: for $t<t'$,
$|C(t;r)-C(t';r)|\le
\sum_j\int|c^r_j(\cdot,t)-c^r_j(\cdot,t')|\dd G_{j,t}
+\big|\sum_j\int c^r_j(\cdot,t')\dd(G_{j,t}-G_{j,t'})\big|$. The
first term is controlled by the uniform time-modulus of the policies
(S4); the second, coupling $G_t$ to $G_{t'}$ at $W_\infty$-distance
$\le\bar s|t-t'|$ (S4) and applying the layer/interior split of
Lemma~\ref{lem:agg}(c) with $\hat\eta$ replaced by $\bar s|t-t'|$,
is $O\big(|t-t'|^{1/2}\big)$ under either density regime ($O(|t-t'|
(1+\log_+))$ under (D1$'$), $O(|t-t'|)$ under (D1), and in the worst
case the $\sqrt{\cdot}$ modulus of consumption near the constraint
gives the square root). Since $Y(\cdot)$ is Lipschitz (autonomous
two-state chain), $\Psi(\Kcal)$ admits the uniform modulus
$\omega(h)=C(\sqrt h+h(1+\log_+(1/h)))$. Arzel\`a--Ascoli gives
relative compactness; Mazur's theorem gives compactness of the closed
convex hull.
\end{proof}

\begin{lemma}[anchoring limit and band invariance]\label{lem:anch}
Let $c_{T,j}:=(u')^{-1}(\partial_a v_{T,j})$ and define
\[
A_0(x)\;:=\;\frac{1}{B}\Big[\textstyle\sum_j
\int_{(\uba,\infty)}c_{T,j}\dd G_{0,j}
+m_1(0)\,\min\{c_{T,1}(\uba),\,y_1+x\,\uba\}-Y(0)\Big],
\qquad x\in\R .
\]
Then: (i) $\sup_{r\in\Kcal}\sup_{t\le T}
|\Psi(r)(t)-A_0(r(t))|=:\varepsilon(T)\to0$ as $T\to0$; (ii) $A_0$ is
nonincreasing and Lipschitz with constant
$\theta_0:=m_1(0)|\uba|/B\le\theta\le\bar\theta<1$, so
$x\mapsto x-A_0(x)$ is strictly increasing with slope at least
$1-\bar\theta$ and $A_0$ has a unique fixed point $r^{**}\in\R$;
(iii) under \textbf{(ANCH)}---$r^{**}\in(r_{\min}+\delta,
r_{\max}-\delta)$ for some $\delta>0$---there is $T_1>0$ such that
$\Psi(\Kcal)\subset\Kcal$ for all $T\le T_1$.
\end{lemma}

\begin{proof}
(i) As $T\to0$: the value functions converge to the terminal data
uniformly on compacts at rate $O(T)$ (the running payoff and the
switching corrections over $[t,T]$ are $O(T)$ by the bounds of
Appendix~\ref{app:household}), so by concavity the interior policies
converge to $c_{T,j}$ locally uniformly; at the constraint the
$\min$-structure of the boundary policy delivers exactly
$\min\{c_{T,1}(\uba),y_1+r(t)\uba\}$ in the limit, consistent with
the terminal boundary layer of Remark~\ref{rem:BClayer}.
Distributions: $W_\infty(G_t,G_0)\le\bar sT\to0$ (S4), and the atom
mass converges to $m_1(0)$ in the binding regime
(Proposition~\ref{prop:atom}). Passing to the limit in
$C(t;r)$ with the split of Lemma~\ref{lem:agg} gives the claim,
uniformly in $r\in\Kcal$ because the limit map depends on the path
only through the one-dimensional value $r(t)$.
(ii) $x\mapsto\min\{c_{T,1}(\uba),y_1+x\uba\}$ is nonincreasing
(slope $0$ or $\uba<0$) with Lipschitz constant $|\uba|$; multiply by
$m_1(0)/B$. Since $A_0$ is bounded above and decreases at most
linearly, $x-A_0(x)$ is strictly increasing (slope
$\ge1-\bar\theta>0$) and changes sign, giving a unique root
$r^{**}$.
(iii) For $x\in[r_{\min},r_{\max}]$,
$|A_0(x)-r^{**}|=|A_0(x)-A_0(r^{**})|\le\bar\theta\,|x-r^{**}|$,
hence
$A_0(x)\ge r^{**}-\bar\theta\,(r^{**}-r_{\min})
=(1-\bar\theta)r^{**}+\bar\theta r_{\min}\ge
r_{\min}+(1-\bar\theta)\delta$, and symmetrically
$A_0(x)\le r_{\max}-(1-\bar\theta)\delta$. Choose $T_1$ with
$\varepsilon(T_1)\le\tfrac12(1-\bar\theta)\delta$; then for $T\le
T_1$ and every $r\in\Kcal$, $\Psi(r)(t)\in[r_{\min},r_{\max}]$ for
all $t$.
\end{proof}

\begin{proof}[Proof of Theorem~\ref{thm:smallT}]
Let $T\le T_1$ (Lemma~\ref{lem:anch}).

\emph{(i) Existence.} Let $X$ be the closed convex hull of
$\Psi(\Kcal)$ in $C([0,T])$: compact and convex by
Lemma~\ref{lem:cpt}, contained in $\Kcal$ by
Lemma~\ref{lem:anch}(iii), and $\Psi(X)\subset
\overline{\mathrm{conv}}\,\Psi(\Kcal)=X$. By Lemma~\ref{lem:agg},
$\|\Psi(r)-\Psi(r')\|_\infty\le\bar\theta\eta+C_1T\eta
+C_2T\eta(1+\log_+(1/\eta))\to0$ as $\eta\to0$, so $\Psi$ is
continuous on $X$. Schauder's fixed-point theorem yields
$r^\ast=\Psi(r^\ast)\in X$. By construction the induced
$(v^{r^\ast},G^{r^\ast},r^\ast)$ satisfies flow clearing
$C=Y+r^\ast B$ at every $t$; with $S(0)=B$,
Lemma~\ref{lem:flow} upgrades this to $S(t)\equiv B$: an equilibrium
transition.

\emph{(ii) Uniqueness under (D1).} By Lemma~\ref{lem:agg}(D1) and
Lemma~\ref{lem:anch}(iii), $\Psi$ is a self-map of the complete
metric space $(\Kcal,\|\cdot\|_\infty)$ with
$\|\Psi(r)-\Psi(r')\|_\infty\le(\bar\theta+C_3T)\,\eta$. Choose
$T^*\le T_1$ with $\bar\theta+C_3T^*\le\tfrac12(1+\bar\theta)<1$;
Banach's fixed-point theorem gives uniqueness in $\Kcal$ and
geometric convergence of the iteration
$r^{n+1}=\Psi(r^n)$ from any $r^0\in\Kcal$---which is precisely the
economic scheme: guess a rate path, solve households backward,
aggregate forward, reset the rate from flow clearing.

\emph{(iii) Ambiguity bound under (D1$'$).} Let $r\ne r'$ be two
equilibria in $\Kcal$ and $\eta=\|r-r'\|_\infty>0$. Both are fixed
points, so Lemma~\ref{lem:agg} gives
$\eta\le\bar\theta\eta+C_1T\eta+C_2T\eta(1+\log_+(1/\eta))$.
Dividing by $\eta>0$,
\[
1\ \le\ \bar\theta+(C_1+C_2)T+C_2T\,\log_+(1/\eta)
\qquad\Longrightarrow\qquad
\eta\ \le\ \exp\!\Big(-\tfrac{1-\bar\theta-(C_1+C_2)T}{C_2T}\Big),
\]
for $T^*$ small enough that the numerator is positive. The same
inequality applied to $r^{n+1}=\Psi(r^n)$ against a fixed point shows
the Picard iterates contract geometrically until they enter a ball of
that exponentially small radius. The estimate does not self-improve
to $\eta=0$: the logarithm grows as $\eta\downarrow0$ on the
favorable side, and the two-directional time coupling---the
anticipation channel propagates backward from $T$, the history
channel forward from $0$---blocks a Gr\"onwall--Osgood closure in
either single time direction.
\end{proof}

\begin{proposition}[diagonal elimination and its price]
\label{prop:diagelim}
Solving the flow condition exactly for $r(t)$ yields
\[
r(t)=\widetilde\Psi(r)(t)
=\frac{C^{\mathrm{int}}(t;r)+m_1(t;r)\,y_1-Y(t)}
{B+m_1(t;r)\,|\uba|},
\]
with $C^{\mathrm{int}}$ the consumption of the unconstrained.
$\widetilde\Psi$ has zero contemporaneous sensitivity and a uniformly
invertible denominator for $B>0$; under (D1),
Theorem~\ref{thm:smallT}(ii) holds for $\widetilde\Psi$ without the
hypothesis $\theta<1$. Under (D1$'$), the separated atom mass is only
$\tfrac12$-H\"older stable in the path
($|m_1(t;r)-m_1(t;r')|\le C\sqrt{T\eta}$, the mismatch mass between
atom and boundary layer), and the loss enters $\widetilde\Psi$ at
order one: no contraction. The trade-off is sharp at the level of
these estimates: aggregates are log-Lipschitz because mass exchanged
between atom and layer moves consumption at second order; any map
separating the atom pays the H\"older price. Under (D1$'$) the
economic threshold $\theta<1$ is genuinely load-bearing.
\end{proposition}

\paragraph{Theorem~\ref{thm:B0atom}.} Under binding on $[0,T]$, atom
outflow is only through switching, giving
$m_1(t)\ge m_1(0)e^{-\lambda_1t}$; the map
$\widetilde\Psi_0(r)(t)=\big(C^{\mathrm{int}}(t;r)+m_1(t;r)y_1
-Y(t)\big)/\big(|\uba|m_1(t;r)\big)$ is then well defined, its
denominator bounded below by $|\uba|m_1(0)e^{-\lambda_1T}$; the
Schauder and contraction steps proceed as in
Theorem~\ref{thm:smallT}, with all constants scaled by
$1/(|\uba|m_1(0))$---whence $T^*$ degrading linearly in $m_1(0)$.
Under (D1$'$), two equilibria satisfy $\eta\le C\sqrt{T\eta}$, i.e.\
$\eta\le CT$; exact uniqueness is open.

\paragraph{Theorem~\ref{thm:mono}.} Standard mixed-monotone iteration
for antitone maps on the order-complete cone of $C([0,T])$:
(H-MONO) makes $\Psi$ antitone, so the two-sided scheme generates
nested order intervals; compactness of $\Psi$'s range upgrades
monotone pointwise convergence to uniform convergence with continuous
limits; any equilibrium is squeezed by induction, and ordered
equilibria coincide by antitonicity. A degree-theoretic continuation
in $T$ from the small-horizon equilibrium (index $+1$) extends
existence along the branch while a priori interiority holds and
$I-D\Psi$ stays invertible; (H-MONO) supplies sign structure but not
spectral control, so this continuation remains conditional---and
Proposition~\ref{prop:multiplicity}(ii) identifies exactly where it
fails for $\gamma>1$.

%======================================================================
\begin{thebibliography}{99}

\bibitem[A\c{c}{\i}kg\"oz(2018)]{Acikgoz2018}
A\c{c}{\i}kg\"oz, \"Omer T. 2018. ``On the Existence and Uniqueness
of Stationary Equilibrium in Bewley Economies with Production.''
\emph{Journal of Economic Theory} 173: 18--55.

\bibitem[Achdou et al.(2014)]{ABLLM2014}
Achdou, Yves, Francisco J. Buera, Jean-Michel Lasry, Pierre-Louis
Lions, and Benjamin Moll. 2014. ``Partial Differential Equation
Models in Macroeconomics.'' \emph{Philosophical Transactions of the
Royal Society A} 372: 20130397.

\bibitem[Achdou et al.(2022)]{AHLLM2022}
Achdou, Yves, Jiequn Han, Jean-Michel Lasry, Pierre-Louis Lions, and
Benjamin Moll. 2022. ``Income and Wealth Distribution in
Macroeconomics: A Continuous-Time Approach.'' \emph{Review of
Economic Studies} 89 (1): 45--86.

\bibitem[Aiyagari(1994)]{Aiyagari1994}
Aiyagari, S. Rao. 1994. ``Uninsured Idiosyncratic Risk and Aggregate
Saving.'' \emph{Quarterly Journal of Economics} 109 (3): 659--684.

\bibitem[Ambrose(2021)]{Ambrose2021}
Ambrose, David M. 2021. ``Existence Theory for a Time-Dependent Mean
Field Games Model of Household Wealth.'' \emph{Applied Mathematics \&
Optimization} 83 (3): 2051--2081.

\bibitem[Auclert et al.(2021)]{AuclertEtAl2021}
Auclert, Adrien, Bence Bard\'oczy, Matthew Rognlie, and Ludwig
Straub. 2021. ``Using the Sequence-Space Jacobian to Solve and
Estimate Heterogeneous-Agent Models.'' \emph{Econometrica} 89 (5):
2375--2408.

\bibitem[Bayer, Rendall, and W\"alde(2019)]{BayerRendallWaelde2019}
Bayer, Christian, Alan D. Rendall, and Klaus W\"alde. 2019. ``The
Invariant Distribution of Wealth and Employment Status in a Small
Open Economy with Precautionary Savings.'' \emph{Journal of
Mathematical Economics} 85: 17--37.

\bibitem[Boppart, Krusell, and Mitman(2018)]{BoppartKrusellMitman2018}
Boppart, Timo, Per Krusell, and Kurt Mitman. 2018. ``Exploiting MIT
Shocks in Heterogeneous-Agent Economies: The Impulse Response as a
Numerical Derivative.'' \emph{Journal of Economic Dynamics and
Control} 89: 68--92.

\bibitem[Bouchut and James(1998)]{BouchutJames1998}
Bouchut, Fran\c{c}ois, and Fran\c{c}ois James. 1998.
``One-Dimensional Transport Equations with Discontinuous
Coefficients.'' \emph{Nonlinear Analysis} 32 (7): 891--933.

\bibitem[Cannarsa and Capuani(2018)]{CannarsaCapuani2018}
Cannarsa, Piermarco, and Rossana Capuani. 2018. ``Existence and
Uniqueness for Mean Field Games with State Constraints.'' In
\emph{PDE Models for Multi-Agent Phenomena}, Springer INdAM Series
28: 49--71.

\bibitem[Cannarsa, Capuani, and Cardaliaguet(2021)]{CCC2021}
Cannarsa, Piermarco, Rossana Capuani, and Pierre Cardaliaguet. 2021.
``Mean Field Games with State Constraints: From Mild to Pointwise
Solutions of the PDE System.'' \emph{Calculus of Variations and
Partial Differential Equations} 60: article 108.

\bibitem[Cannarsa and Sinestrari(2004)]{CannarsaSinestrari2004}
Cannarsa, Piermarco, and Carlo Sinestrari. 2004. \emph{Semiconcave
Functions, Hamilton--Jacobi Equations, and Optimal Control}.
Boston: Birkh\"auser.

\bibitem[Capuzzo-Dolcetta and Lions(1990)]{CDL1990}
Capuzzo-Dolcetta, Italo, and Pierre-Louis Lions. 1990.
``Hamilton--Jacobi Equations with State Constraints.''
\emph{Transactions of the American Mathematical Society} 318 (2):
643--683.

\bibitem[Cardaliaguet and Lehalle(2018)]{CardaliaguetLehalle2018}
Cardaliaguet, Pierre, and Charles-Albert Lehalle. 2018. ``Mean Field
Game of Controls and an Application to Trade Crowding.''
\emph{Mathematics and Financial Economics} 12: 335--363.

\bibitem[Chen et al.(2025)]{ChenEtAl2025}
Chen, Cangxiong, Sigmund Ellingsrud, Fabian Harang, Alfonso
Irarrazabal, and Avi Mayorcas. 2025. ``Stochastic Analysis of
Overlapping Generations Models under Incomplete Markets.''
arXiv:2509.05170.

\bibitem[Cirant, Gianni, and Mannucci(2020)]{CirantGianniMannucci2020}
Cirant, Marco, Roberto Gianni, and Paola Mannucci. 2020. ``Short-Time
Existence for a General Backward--Forward Parabolic System Arising
from Mean-Field Games.'' \emph{Dynamic Games and Applications} 10
(1): 100--119.

\bibitem[Crandall, Ishii, and Lions(1992)]{CIL1992}
Crandall, Michael G., Hitoshi Ishii, and Pierre-Louis Lions. 1992.
``User's Guide to Viscosity Solutions of Second Order Partial
Differential Equations.'' \emph{Bulletin of the American Mathematical
Society} 27 (1): 1--67.

\bibitem[Engler and Lenhart(1991)]{EnglerLenhart1991}
Engler, Hans, and Suzanne M. Lenhart. 1991. ``Viscosity Solutions for
Weakly Coupled Systems of Hamilton--Jacobi Equations.''
\emph{Proceedings of the London Mathematical Society} 63 (1):
212--240.

\bibitem[Filippov(1960)]{Filippov1960}
Filippov, Aleksei F. 1960. ``Differential Equations with
Discontinuous Right-Hand Side.'' \emph{Matematicheskii Sbornik} 51:
99--128. English translation: \emph{AMS Translations, Series 2} 42
(1964): 199--231.

\bibitem[Fujii and Takahashi(2022)]{FujiiTakahashi2022}
Fujii, Masaaki, and Akihiko Takahashi. 2022. ``A Mean Field Game
Approach to Equilibrium Pricing with Market Clearing Condition.''
\emph{SIAM Journal on Control and Optimization} 60 (1): 259--279.

\bibitem[Gomes and Sa\'ude(2021)]{GomesSaude2021}
Gomes, Diogo A., and Jo\~ao Sa\'ude. 2021. ``A Mean-Field Game
Approach to Price Formation.'' \emph{Dynamic Games and Applications}
11 (1): 29--53.

\bibitem[Graber and Bensoussan(2018)]{GraberBensoussan2018}
Graber, P. Jameson, and Alain Bensoussan. 2018. ``Existence and
Uniqueness of Solutions for Bertrand and Cournot Mean Field Games.''
\emph{Applied Mathematics \& Optimization} 77 (1): 47--71.

\bibitem[Graber and Mayorga(2021)]{GraberMayorga2021}
Graber, P. Jameson, and Sergio Mayorga. 2021. ``A Note on
Deterministic Mean Field Games of Controls with State Constraints:
Existence of Mild Solutions.'' arXiv:2109.11655 (revised 2023).

\bibitem[H\"ofer(2025)]{Hofer2025}
H\"ofer, F. 2025. ``Stationary Heterogeneous-Agent Models in
Continuous Time.'' arXiv:2510.26065 (revised 2026).

\bibitem[Huggett(1993)]{Huggett1993}
Huggett, Mark. 1993. ``The Risk-Free Rate in Heterogeneous-Agent
Incomplete-Insurance Economies.'' \emph{Journal of Economic Dynamics
and Control} 17 (5--6): 953--969.

\bibitem[Ishii and Koike(1991)]{IshiiKoike1991}
Ishii, Hitoshi, and Shigeaki Koike. 1991. ``Viscosity Solutions for
Monotone Systems of Second-Order Elliptic PDEs.''
\emph{Communications in Partial Differential Equations} 16 (6--7):
1095--1128.

\bibitem[Kaplan, Moll, and Violante(2018)]{KMV2018}
Kaplan, Greg, Benjamin Moll, and Giovanni L. Violante. 2018.
``Monetary Policy According to HANK.'' \emph{American Economic
Review} 108 (3): 697--743.

\bibitem[Kobeissi(2022)]{Kobeissi2022}
Kobeissi, Ziad. 2022. ``Mean Field Games with Monotonous Interactions
through the Law of States and Controls of the Agents.'' \emph{NoDEA
Nonlinear Differential Equations and Applications} 29: article 52.

\bibitem[Porretta(2014)]{Porretta2014}
Porretta, Alessio. 2014. ``On the Planning Problem for the Mean Field
Games System.'' \emph{Dynamic Games and Applications} 4 (2):
231--256.

\bibitem[Poupaud and Rascle(1997)]{PoupaudRascle1997}
Poupaud, Fr\'ed\'eric, and Michel Rascle. 1997. ``Measure Solutions
to the Linear Multi-Dimensional Transport Equation with Non-Smooth
Coefficients.'' \emph{Communications in Partial Differential
Equations} 22 (1--2): 337--358.

\bibitem[Shigeta(2023)]{Shigeta2023}
Shigeta, Yuki. 2023. ``Existence of Invariant Measure and Stationary
Equilibrium in a Continuous-Time One-Asset Aiyagari Model: A Case of
Regular Controls under Markov Chain Uncertainty.'' SSRN Working Paper
4510345; companion paper SSRN 4510320.

\bibitem[Soner(1986a)]{Soner1986a}
Soner, Halil Mete. 1986a. ``Optimal Control with State-Space
Constraint I.'' \emph{SIAM Journal on Control and Optimization} 24
(3): 552--561.

\bibitem[Soner(1986b)]{Soner1986b}
Soner, Halil Mete. 1986b. ``Optimal Control with State-Space
Constraint II.'' \emph{SIAM Journal on Control and Optimization} 24
(6): 1110--1122.

\end{thebibliography}

\end{document}
